
\section*{Introduction}
\addcontentsline{toc}{section}{\protect\numberline{}Introduction}

    The \textbf{relevance} of this work is determined by the pivotal role that Large Language Models (LLMs) currently play in various domains including machine learning \cite{machine_learning}, computer vision \cite{computer_vision}, and natural language processing \cite{brown2020languagemodelsfewshotlearners}. These pre-trained LLMs are actively employed in numerous industrial applications, prized for their proficiency in tasks such as in-context learning \cite{brown2020languagemodelsfewshotlearners}, instruction following \cite{instruction_following}, and structured reasoning \cite{reasoning}. However, the capabilities of a pre-trained LLM may not be sufficient to tackle
    more specialized tasks, for example in medical or scientific domains. This necessitates extensive fine-tuning to tailor the models for specific downstream tasks. However, the fine-tuning process is extremely resource-consuming and demands signi-ficant
    memory and time resources resulting in low training throughput \cite{malladi2024finetuninglanguagemodelsjust}. This way, the evolvement of novel fine-tuning strategies that reduce time and memory costs without sacrificing model quality is of great importance.

    At the same time, given the vast parameter space of modern LLMs, which results in significant over-parameterization \cite{overparametrization}, identifying and leveraging the most essential model structures becomes crucial. These structures are vital as they disproportionately impact model performance \cite{yu2024superweightlargelanguage}. Effective utilisation of LLMs in various scenarios -- whether pre-training, fine-tuning, or inference -- demands a nuanced understanding of "importance" to discern these critical structures. This requirement extends across multiple optimization strategies including Neural Architecture Search \cite{LPZero} (NAS), \cite{synflow}, \cite{snip}, Parameter-Efficient Fine-Tuning (\peft{}) \cite{adalora}, \cite{lisa}, quantization \cite{quantization_1}, \cite{int8}, \cite{SqueezeLLM}, sparsification \cite{chess}, and model pruning \cite{uidl}, \cite{sleb}, \cite{wanda}. Each of these strategies aims to optimize different aspects of model architecture, yet they often rely on disparate heuristics for estimating importance, many of which lack general applicability or direct correlation with model quality \cite{LPZero}.

    This division highlights a glaring gap: the absence of a general methodology capable of assessing the importance of substructures across different \llm{} applications. Addressing this gap is not only crucial for enhancing model efficiency but also for advancing the methodologies used in model optimization, in particular, in the field of model fine-tuning acceleration. 

    The \textbf{practical significance} of this research is that it introduces a comprehensive overview of various existing importance estimation approaches as applied in distinct scenarios and presents a thorough examination of the most prominent approaches for importance analysis of Transformer blocks through the example of fine-tuning acceleration.

    The \textbf{novelty} of this work can be summarized as follows:
    \begin{itemize}
        \item This work introduces a comprehensive overview of existing methodologies for importance estimation of different structures, such as separate weights, layers, Transformer blocks, and even predicted architectures within the field of \llm{} across diverse application contexts. 
        \item This research provides evaluation of the most prominent approaches for esti-mating Transformer blocks importance within the context of accelerating the fine-tuning of \llms{}.
    \end{itemize}

    The \textbf{goal} of the present study is to accelerate the fine-tuning process by introducing the importance analysis of Transformer blocks within the \llm{}.

    To achieve the goal, the following \textbf{objectives} must be accomplished: 
    \begin{enumerate}
        \item Conduct comprehensive background study on the topic of fine-tuning \\ acceleration directions. 
        \item Provide an overview of existing methodologies for estimation the saliency of various structures within \llm{} in diverse application scenarios. 
        \item Conduct evaluation of the existing importance estimation techniques for \\ Transformer blocks in the context of accelerating the fine-tuning of \llm{}.
        \item Analyze obtained results and draw conclusion.
    \end{enumerate}

    The study employs a combination of \textbf{methodological approaches} to adress the research objectives. Firstly, it adopts a method of theoretical analysis and synthesis of existing literature to establish a solid theoretical foundation for the research. Secondly, programming methods are utilized to analyze the importance of \llms{}' Transdormer blcosk. Lastly, the study employs Deep Learning techniques to fine-tune \llms{}.

    The \textbf{object} of this study is the importance analysis of different structures within the \llm{}. The \textbf{subject} of the study is a thorough examination of the Transformer blocks-based importance approaches through the acceleration of fine-tuning process.

    This study was conducted as part of a \textbf{research project at the Huawei} dedicated to the fine-tuning acceleration of \llms{}.

